\section{Tracing Frameworks}

Tracing frameworks provide the internal infrastructure required to interface
event sources and front ends. The Linux kernel provides multiple frameworks
tailored to different use cases. A framework should support the registration of
event sources, which can be categorized as either static or dynamic, to capture
their side effects when such events are triggered. User-space interacts with a
tracing frameworks through a well-defined stable user-space API, such as through
a filesystem or a dedicated system call. A front-end provides a more convenient
API to the user for managing event registration, de-registration and consuming
its output.

\subsection{Ftrace}

Ftrace is used to support the registration, deregistration, reading, writing and
other operations on various types of trace events. It provides a \textit{Tracer}
abstraction to manage a set of events together, and a data structure to store
and interact with all the traced data. It also provides other infrastructure for
specific tracers, such as hooks to any kernel function and code patching
routines.

\textit{Ftrace} was originally implemented in the Linux kernel to support
function tracing, that is, capturing information such as arguments and function
name whenever a kernel function is entered or exited in real time. Overtime the
same infrastructure was reused for other use cases, therefore the name "ftrace"
does not retain any of its original meaning. For this reason, we will use the
term "function tracer" to refer to the specific tracer used for functions, and
"ftrace" for the wider infrastructure.

\subsubsection*{Global Ring Buffer}

A core component of ftrace is its \textit{global ring buffer}, represented on
its highest abstraction level by the \verb|struct trace_array|, which provides
an interface to write an event's data and consume it. It supports both a
producer-consumer mode and an overwrite mode. The two modes are distinguished by
how writers behaves when the buffer is full: in such case, if the buffer is in
prodicer-consimer mode, then the producer fails to write, otherwise in overwrite
mode it overwrite the data.

The design of the ring buffer had to met specific requirements because trace
events are expected to be originated from any context, unless specifically
disabled such on the tracing infrastructure itself to avoid recursion, and at
high frequency.

Core requirements are, but not limited to, the following:

\begin{itemize}

  \item writers should be able to write from interruptable contexts with success
    when interrupted.

  \item many writers and readers are expected, in multiple CPU cores.

  \item the synchronization overhead must be close to zero due to
the potential high amount of concurrent writes from high frequency events.

\end{itemize}

The solution implemented in Linux in the file
\verb|linux/kernel/trace/ring_buffer.c| honors all these requirements using a
lockless design, for both reader-to-writer and writer-to-writer interactions.

Some key design decisions are now briefly discussed.

Each CPU uses its own per-CPU ringbuffer containing a doubly-linked list of
pages; therefore writers on different CPUs do not need to synchronize among each
other. In order to read a full trace in chronological order, a reader iterates
over all the buffers using timestamp information to determine the order of
events.

To support preemption, writers use a \textit{reserve and commit} pattern where
the writer first reserves allocation space in a per-CPU ring buffer, then it
writes with exclusive access to that region, and finally commits the region. The
commit will be in a pending state until all regions before it are also
committed; therefore when a writer reserves a region but gets preempted by
another writer, then the new writer will not have access to the reserved region
and its changes will not get registered until the original writer completes.
This mechanism ensures that a reader will advance only if all preempted writes
are committed.

In C code, the commit pattern looks like the following function from
\verb|linux/kernel/trace/trace_uprobe.c|, commented and simplified for
readability:

\begin{verbatim}
static void __uprobe_trace_func(struct trace_uprobe *tu,
				unsigned long func, struct pt_regs *regs,
				struct uprobe_cpu_buffer *ucb,
				struct trace_event_file *trace_file)
{
  struct trace_event_buffer fbuffer;
  void *data;
  int size = ...;

  /* First, reserve an area */
  entry = trace_event_buffer_reserve(&fbuffer, trace_file, size);
  if (!entry)
    return;

  data = /* get the data pointer from and entry */

  /* Second, copy the data into the ring buffer page */
  memcpy(data, ucb->buf, ucb->dsize);

  /* Third, commit */
  trace_event_buffer_commit(&fbuffer);
}
\end{verbatim}

Another key detail is to have the reader page fully extern from the ring buffer;
therefore, writers never interact with the reader page and readers swap the
reader page via atomic operations so no lock is used on the same CPU ring
buffer.

Because string formatting is expensive, the event is stored in the ring buffer
in binary form without any formatting for the end-user. When an end-user reads
the ring buffer though the user-space API, the relevant contents of the ring
buffer are copied to a temporary buffer (\verb|struct trace_seq|) where they are
formatted with a printf-style function, \verb|trace_seq_vprintf|, and copied to
the user's memory. Moving the formatting procedure to the reader improves the
writer's latency which is critical for high frequency events.

The user can read the formatted content of all ring buffers by reading the file
\textit{\$TRACEFS\_ROOT / trace} without consuming the buffer, or
\textit{\$TRACEFS\_ROOT / trace\_pipe} to consume it. Additionally, he can get
information about the buffer state from several files beginning with
\verb|buffer_*|. To access a particular per-CPU buffer, he can use the files
from the directory \textit{\$TRACEFS\_ROOT / per\_cpu / \$CPU\_NUMBER /}.

\subsubsection*{The tracer struct}

An important abstraction in ftrace is the \textit{Tracer} represented as the
\verb|struct tracer|. This is meant to provide a way for the end-user to enable
or disable a set of trace events with specific formatting. Any tracer implements
a set of functions to be invoked during its lifetime, such as initialization,
starting, stopping, reading, etc. An use case, for example, would be to trace
all the functions executed, or to trace all context switches. The user interacts
with the tracer, such as to change its execution state or output options, via
the \textit{tracefs} filesystem API. Many tracers are available, but only one
tracer can be enabled at a time.

The default tracer is the \textit{nop} tracer which is a "dummy" tracer that does
not enable any event. The struct instantiation would look like this:

\begin{verbatim}
/* Supported output options */
enum {
	TRACE_NOP_OPT_ACCEPT = 0x1,
	TRACE_NOP_OPT_REFUSE = 0x2
};

/* Options for the tracer */
static struct tracer_opt nop_opts[] = {
	{ TRACER_OPT(test_nop_accept, TRACE_NOP_OPT_ACCEPT) },
	{ TRACER_OPT(test_nop_refuse, TRACE_NOP_OPT_REFUSE) },
	{ } /* Always set a last empty entry */
};

static struct tracer_flags nop_flags = {
	.val = 0, /* By default: all flags disabled */
	.opts = nop_opts
};

struct tracer nop_trace __read_mostly =
{
	.name		= "nop",
	.init		= nop_trace_init,  /* dummy functions, do not do anything */
	.reset		= nop_trace_reset,
	.flags		= &nop_flags,
	.set_flag	= nop_set_flag,
	.allow_instances = true,
};
\end{verbatim}

Where the \verb|init| and \verb|reset| members are functions, which in this case
do nothing. The tracer is registered in the ftrace infrastructure via
the \verb|register_tracer| function which is called during kernel startup.

The following is a non exhaustive list of available tracers:

\begin{itemize}

\item \textit{function}: trace all kernel functions when entered.

\item \textit{function\_graph}: trace both function entry and exit. It can be
used to draw a graph of function calls.

\item \textit{blk}: block tracer, for detailed IO tracing.

\item \textit{hwlat}: hardware latency tracer.

\item \textit{irqsoff}: trace the areas that disable interrupts and saves the
trace with the longest maximum latency.

\item \textit{preemptoff}: similar to \textit{irqsoff} but traces and records
the amount of time for which preemption is disabled.

\item \textit{wakeup}: traces and records the maximum latency that it takes for
the highest priority task to get scheduled after it has been woken up.

\end{itemize}

The user can query a list of the tracers supported by its running kernel by
reading the file \textit{\$TRACEFS\_ROOT / available\_tracers}. He can select a
tracer by writing its name to the file \textit{\$TRAEFS\_ROOT / current\_tracer}
and tune its options from \textit{\$TRACEFS\_ROOT / options / \$OPTION\_NAME}.

\subsubsection*{Static function tracing}

To support tracing functions enter and exit events, ftrace implements a
system to instrument any function in order to alter the execution and log the
event. For clarity, in the following section we will refer as \textit{trace}
function or \textit{tracing} function to the function that performs the writing
of the event onto the ftrace ringbuffer, and as the \textit{traced} function to
the function that is being instrumented.

Whenever a C function is executed, meaning that the program counter reaches the
first instruction belonging the function, the execution should jump to a tracing
function and returns to the traced function without altering the registers.
Likewise, when the program counter reaches the end of the function and returns,
it should jump to a tracing function and later return to the original return
address. This behaviour should be supported transparently, without altering
existing code or to expect developers to put instrumentation jumps onto any
function.

This is implemented by the \textit{gcc} compiler via the \verb|-pg| command line
flags. When these flags are enabled, the compiler will generate a jump to a
trace function called \verb|mcount| after the prologue of any function on ARM,
or a jump to \verb|__fentry__| before the prologue on the x86-64 architecture.
For function exits, the trace function will change the return address of the
traced function with the address of another trace function, which is saved on
the stack as of C calling conventions. This is referred to as "static function
tracing" because the compiler instruments each function permanently at compile
time.

For example, the following \verb|sum| function would be compiled by gcc version
16.1 into x86-64 assembly with \verb|-O0| flag as shown below. On the left the
code in C, on the right its compiled code in assembly.

\noindent
\vspace{0.5em}

\begin{minipage}[t]{0.47\textwidth}
\begin{verbatim}
int sum(int a, int b) {
    return a + b;
}
\end{verbatim}
\end{minipage}%
\hfill
\vrule % This creates the vertical dividing bar
\hfill
\begin{minipage}[t]{0.47\textwidth}
\begin{verbatim}
"sum(int, int)":
        lea     eax, [rdi+rsi]
        ret
\end{verbatim}
\end{minipage}

\noindent
\vspace{0.5em}

While the same compiler would compile that same function as follows with the
\verb|-pg| flag:

\vspace{0.5em}

\begin{minipage}[t]{0.47\textwidth}
\begin{verbatim}
int sum(int a, int b) {
    return a + b;
}
\end{verbatim}
\end{minipage}%
\hfill
\vrule % This creates the vertical dividing bar
\hfill
\begin{minipage}[t]{0.47\textwidth}
\begin{verbatim}
"sum(int, int)":
        call    __fentry__
        lea     eax, [rdi+rsi]
        ret
\end{verbatim}
\end{minipage}

\vspace{0.5em}

A shortcoming of this implementation is that each function suffers the cost of a
jump and a return at minimum, causing some performance loss in case tracing is
not turned on. This issue is solved with "dynamic function tracing", which we
will discuss next.

Static function tracing can be enabled via the \verb|CONFIG_FUNCTION_TRACER|
compile option.

\subsubsection*{Dynamic function tracing}

In order not to have any performance penalty when tracing is supported but
disabled, ftrace provides a way to dynamically patch functions to add or remove
tracing instrumentation. This is enabled with the compile option
\verb|DYNAMIC_FTRACE|.

If dynamic function tracing is enabled, the kernel will patch all calls to
\verb|mcount| or \verb|__fentry__| with \verb|NOP| operations at boot-time. When
tracing is enabled, the same \verb|NOP| instruction are patched to a call to
\verb|ftrace_caller| or \verb|ftrace_regs_caller|. This code patching allows
for no overhead in case tracing is disabled.

The information of all the locations of all \verb|mcount| calls is generated at
compile time via the \verb|recordmcount| program which is invoked by the build
infrastructure. The tool stores all the locations of these functions in the
\verb|__mcount_loc| ELF section, which is linked in the original object file.

The tracing function \verb|ftrace_caller| is also referred to as a "trampoline"
because its instructions are dynamically patched to jump to the functions which
are currently enabled depending on the current tracer o particular events,
represented in the \verb|struct trace_ops|. Ftrace provides the function
\verb|register_ftrace_function| to register a trampoline which is created by an
architecture dependent function, \verb|create_trampoline|.

For example, consider the following assembly function:

\begin{verbatim}
some_kernel_func:
  PUSH RBP
  MOV RSP, RBP
  ...
  RET
\end{verbatim}

When tracing disabled but dynamic tracing is supported, some \verb|NOP|
operations would be patched in:

\begin{verbatim}
# Symbol where instructions like: CALL some_kernel_func will jump to.
some_kernel_func:
  NOP
  NOP
# actual start of the function
  PUSH RBP
  MOV RSP, RBP
  ...
  RET
\end{verbatim}

When tracing is enabled, these same \verb|NOP| functions are replaced with
\verb|ftrace_caller|:

\begin{verbatim}
# Symbol where instructions like: CALL some_kernel_func will jump to.
some_kernel_func:
  CALL ftrace_caller
  NOP
# actual start of the function
  PUSH RBP
  MOV RSP, RBP
  [...]
  RET
\end{verbatim}

\subsection{Perf}

TODO


\subsection{Others}

There are some vendor specific systems and audit.

There are some out of tree systems:

  - systemtap
  - LTTng
  - ktap
  - dtrace4linux
  - DTrace
  - sysdig
  - ...

TODO
